% Slide 7 — Lacuna na literatura
\begin{frame}{A lacuna não é falta de algoritmos}
  \begin{columns}[T,onlytextwidth]
    \column{0.55\textwidth}
      \textbf{O que já existe}
      \begin{itemize}
        \small
        \item Modelos analíticos clássicos (EOQ, $(s,S)$, Jornaleiro)
        \item Políticas parametrizadas por meta-heurísticas
        \item Políticas aprendidas por reforço
        \item Comparações isoladas entre famílias
      \end{itemize}
    \column{0.45\textwidth}
      \textbf{O que falta}
      \begin{itemize}
        \small
        \item \highlight{Seleção contextual} validada em dados reais
        \item Evidência em \highlight{varejo brasileiro} intermitente
        \item Ligação explícita entre perfil e política dominante
      \end{itemize}
  \end{columns}
  \vfill
  \begin{alertblock}{Posicionamento}
    \small O problema se aproxima da \textbf{seleção de algoritmos}
    (Rice, 1976) -- escolher, para cada instância, a alternativa com melhor
    desempenho esperado.
  \end{alertblock}
  \note{
    A literatura de controle de inventário já oferece muitas alternativas:
    modelos analíticos, meta-heurísticas, aprendizado por reforço. O que
    falta não é mais um algoritmo -- é um mecanismo de seleção contextual
    validado sobre dados reais, especialmente em um contexto pouco coberto,
    que é o varejo brasileiro com demanda intermitente. A formulação que eu
    adoto se aproxima do problema clássico de seleção de algoritmos, de
    Rice, 1976: dado um conjunto de instâncias e um portfólio de
    alternativas, aprender a escolher a alternativa mais adequada para
    cada instância.
  }
\end{frame}
